
\documentclass[letter, 11pt]{article}

    %%%%%%%%%%%%%%%%%%%%%%%%%%%%%%
    \oddsidemargin  0.0in
    \evensidemargin 0.0in
    \textwidth      6.6in
    \headheight     0.0in
    \topmargin      0.0in
    \textheight=9.0in
    \parindent=0in
    %%%%%%%%%%%%%%%%%%%%%%%%%%%%%%

\include{MathCommands}

\usepackage{amsthm,amsfonts,amssymb,amsmath}
\usepackage{graphicx}   %To include graphics
\usepackage[small,compact]{titlesec}

\usepackage[utf8]{inputenc}
\usepackage{bm}
\usepackage[francais]{babel}

\usepackage{fancyhdr}
\pagestyle{fancy}

\renewcommand{\headrulewidth}{1pt}
\fancyhead[C]{Université Laval}
\fancyhead[L]{STT-2100 Régression}
\fancyhead[R]{A.-S. Charest/E. Reny-Nolin}


\begin{document}

\vspace*{0.05 in}
\begin{center}
{\Large Exercices \\
Chapitre 5 - Ajuster un modèle de RLM}
\end{center}
\vspace*{0.1 in}

\begin{enumerate}\itemsep1cm


%%%%%%%%%%%%%%%%%%%%%%%%%%%%%%%%%%%%%%%%%%%%%%%%%%%%%%%%%%%%%%%%%%%%%%%%%%%%%%%%%%%%%%%%%%
\item Marie a 34 ans,  16 ans de scolarité, et elle gagne $71\,000$ \$ par année. Marc est allé à l'école pendant 21 ans, il gagne $56\,000$ \$, et vient de souffler ses 30 bougies. André, 50 ans, a un revenu annuel de $91\,000$ \$ et 16 ans de scolarité. Annie, 40 ans, a 14 ans de scolarité et gagne $45\,000$ \$.
    
    On veut utiliser un modèle linéaire multiple pour prédire le revenu à partir de l'âge et de la scolarité. 
    
    \begin{itemize}
    \item[(a)] Écrivez le modèle sous sa forme matricielle, et préciser explicitement le contenu de chaque vecteur ou matrice dont les termes sont connus (avant estimation). 
    \item[(b)] Quel est le rang de la matrice $\bm{X}$?
    \end{itemize}

%%%%%%%%%%%%%%%%%%%%%%%%%%%%%%%%%%%%%%%%%%%%%%%%%%%%%%%%%%%%%%%%%%%%%%%%%%%%%%%%%%%%%%%%%%
\item Soit la matrice $\bm{A}=\left[\begin{array}{cc}3&1\\4&2\end{array}\right]$ et le vecteur $\bm{v}'=[v_1\;v_2]$.

\bigskip
    \begin{itemize}\itemsep4mm
    \item[(a)] Déterminez $\dfrac{\partial \bm{v'A}}{\partial \bm{v}}$.
    \item[(b)] Déterminez $\dfrac{\partial \bm{v'Av}}{\partial \bm{v}}$.
    \end{itemize}


%%%%%%%%%%%%%%%%%%%%%%%%%%%%%%%%%%%%%%%%%%%%%%%%%%%%%%%%%%%%%%%%%%%%%%%%%%%%%%%%%%%%%%%%%%
\item Démontrez que la matrice chapeau\\
    \begin{itemize}
    \item[(a)] est idempotente\\
    \item[(b)] est symétrique \\
    \item[(c)] a une trace égale à $p' = p+1$. (Utilisez le fait que  tr$(AB)$ = tr$(BA)$ si les dimensions le permettent.)

    \end{itemize}





\end{enumerate}

\end{document} 